% =============================================================================
% Capítulo 3: Arquitectura Propuesta (nombre literal del enunciado)
% Las versiones de la Tabla de herramientas se verificaron en el equipo de
% desarrollo el 05-11-2026 (ver docs/bitacora.md).
% =============================================================================
\section{Arquitectura Propuesta}
\label{sec:arquitectura}

La solución sigue una arquitectura por capas, desde las fuentes de datos hasta
el usuario final (Figura~\ref{fig:arquitectura}). Los datos se extraen de las
fuentes públicas con Python, se integran con procesos ETL en Apache Hop y se
almacenan en un Datawarehouse en PostgreSQL. Desde ahí, un cubo semántico en
Cube define las dimensiones y métricas que consume el dashboard en Streamlit,
dirigido a la gerencia comercial y de marketing. Todos los servicios corren en
contenedores Docker, por lo que el entorno se reproduce igual en Linux y en
Windows.

\begin{figure}[H]
  \centering
  \caption{Arquitectura de la solución BI por capas}
  \label{fig:arquitectura}
  \contorno[6pt]{%
  \begin{tikzpicture}[
      font=\linespread{1}\scriptsize,
      etiqueta/.style={fill=azulutem, text=white, font=\linespread{1}\scriptsize\bfseries,
                       align=flush center, text width=2.3cm, minimum height=1.15cm,
                       inner sep=2pt},
      comp/.style={draw=azulutem, thick, rounded corners=2pt, fill=azulutem!6,
                   align=flush center, inner sep=3pt, minimum height=0.95cm},
      fuente/.style={comp, draw=gray!70, fill=gray!8, text width=1.55cm,
                     inner sep=1.5pt, minimum height=1.05cm},
      hecho/.style={comp, draw=verdeutem!75!black, fill=verdeutem!18},
      flecha/.style={-{Stealth[length=2mm]}, thick, draw=gray!65!black},
      nota/.style={font=\linespread{1}\tiny\itshape, text=gray!40!black}]

    % Etiquetas de las capas
    \node[etiqueta] at (1.25, 0)     {Fuentes\\de datos};
    \node[etiqueta] at (1.25,-1.75)  {Extracción};
    \node[etiqueta] at (1.25,-3.5)   {Integración\\(ETL)};
    \node[etiqueta] at (1.25,-5.25)  {Almacena-\\miento};
    \node[etiqueta] at (1.25,-7.0)   {Explotación};
    \node[etiqueta] at (1.25,-8.75)  {Usuario};

    % Capa 1: una fuente por línea de negocio y las fuentes de apoyo
    \foreach \nombre/\texto [count=\i] in {
        f1/\textbf{Buses}\\ANTT,
        f2/\textbf{Trenes}\\UK Train Rides,
        f3/\textbf{Taxis}\\NYC TLC,
        f4/\textbf{Logística}\\Olist,
        f5/\textbf{Aerolíneas}\\BTS DB1B,
        f6/\textbf{Delivery}\\Delivery Center,
        f7/{\textbf{Apoyo}\\Censo 2024, BCCh}} {
      \node[fuente] (\nombre) at ({3.85+(\i-1)*1.84}, 0) {\texto};
    }

    % Capa 2: extracción
    \node[comp, text width=12.4cm] (ext) at (9.35,-1.75)
      {\textbf{Python} (pandas, SQLAlchemy, psycopg): descarga de las fuentes y
       archivos originales en \texttt{data/raw/}};

    % Capa 3: integración
    \node[comp, text width=12.4cm] (hop) at (9.35,-3.5)
      {\textbf{Apache Hop} (pipelines y workflows, ejecutados con
       \texttt{hop-run}): limpieza y validación; normalización de zonas,
       monedas (CLP), unidades y catálogos; asignación de clientes sintéticos;
       carga de dimensiones y hechos; estadísticas de procesamiento};

    % Capa 4: almacenamiento en PostgreSQL 17
    \node[comp, text width=3.9cm] (stg) at (5.05,-5.25)
      {\textbf{PostgreSQL 17}\\esquema \texttt{stg} (staging)};
    \node[hecho, text width=3.9cm] (dwh) at (10.0,-5.25)
      {\textbf{PostgreSQL 17}\\esquema \texttt{dwh} (modelo estrella)};
    \node[hecho, dashed, text width=2.9cm] (pga) at (14.3,-5.25)
      {\textbf{pgAdmin 4}\\administración};

    % Capa 5: explotación
    \node[comp, text width=4.4cm] (cube) at (9.6,-7.0)
      {\textbf{Cube}\\cubo semántico:\\dimensiones, medidas y KPI};
    \node[comp, text width=2.9cm] (st) at (14.3,-7.0)
      {\textbf{Streamlit + Plotly}\\dashboard};

    % Capa 6: usuario
    \node[comp, draw=verdeutem!75!black, fill=white, text width=12.4cm] (us) at (9.35,-8.75)
      {\textbf{Gerencia comercial y de marketing} (perfil táctico): segmentación
       de clientes, campañas, descuentos y ofertas entre líneas};

    % Flujos
    \foreach \i in {1,...,7} {
      \draw[flecha] (f\i.south) -- (f\i.south |- ext.north);
    }
    \draw[flecha] (ext.south) -- (ext.south |- hop.north);
    \draw[{Stealth[length=2mm]}-{Stealth[length=2mm]}, thick, draw=gray!65!black]
      (stg.north |- hop.south) -- node[nota, right] {carga y lectura} (stg.north);
    \draw[flecha] (dwh.north |- hop.south) -- node[nota, right] {carga} (dwh.north);
    \draw[thick, dashed, draw=gray!65!black] (dwh.east) -- (pga.west);
    \draw[flecha] (dwh.south) -- node[nota, right] {consultas SQL} (dwh.south |- cube.north);
    \draw[flecha] (cube.east) -- node[nota, above] {API} (st.west);
    \draw[flecha] (st.south) -- (st.south |- us.north);

    % Leyenda
    \tikzset{muestra/.style={thick, rounded corners=1pt, minimum height=0.35cm,
                             minimum width=0.6cm, inner sep=0pt}}
    % Cada elemento parte donde termina el texto anterior, para que la leyenda
    % no se monte sobre sí misma con ninguna tipografía
    \node[muestra, draw=verdeutem!75!black, fill=verdeutem!18, anchor=west]
      (l1) at (0.4,-10.0) {};
    \node[anchor=west] (t1) at (l1.east) {Implementado en el Avance N°1};
    \node[muestra, draw=azulutem, fill=azulutem!6, anchor=west]
      (l2) at ([xshift=6mm]t1.east) {};
    \node[anchor=west] (t2) at (l2.east) {Se implementa en los Avances N°2 y N°3};
    \node[muestra, draw=gray!70, fill=gray!8, anchor=west]
      (l3) at ([xshift=6mm]t2.east) {};
    \node[anchor=west] at (l3.east) {Fuente externa};
  \end{tikzpicture}}
  \nota{Elaboración propia. BCCh: Banco Central de Chile (tipos de cambio).
  Todos los servicios se ejecutan en contenedores Docker Compose.}
\end{figure}

\subsection{Componentes y Herramientas}

La Tabla~\ref{tab:herramientas} detalla las herramientas de cada capa. Las
versiones son las instaladas en el equipo de desarrollo y quedan fijadas en
\bd{docker-compose.yml} (imágenes de los servicios) y en
\bd{requirements.txt} (paquetes de Python).

\begin{table}[H]
  \centering
  \caption{Herramientas de la arquitectura por capa}
  \label{tab:herramientas}
  \small
  \begin{tabularx}{\linewidth}{@{}>{\raggedright\arraybackslash}p{2.9cm}
                                  >{\raggedright\arraybackslash}p{3.6cm}
                                  >{\raggedright\arraybackslash}p{2.3cm}
                                  >{\raggedright\arraybackslash}X@{}}
    \toprule
    \textbf{Capa} & \textbf{Herramienta} & \textbf{Versión} & \textbf{Función} \\
    \midrule
    Extracción & Python & 3.13.5 & Lenguaje de los scripts de extracción y de apoyo \\
               & pandas & 3.1.0 & Lectura y preparación de archivos CSV, Excel y Parquet \\
               & SQLAlchemy y psycopg & 2.1.6 y 3.3.8 & Conexión de Python con PostgreSQL \\
    \midrule
    Integración & Apache Hop~\cite{apache_hop} & 2.20.0 & Pipelines y workflows ETL, ejecutados por línea de comandos (\bd{hop-run}); requiere Java 17 o superior \\
    \midrule
    Almacenamiento & PostgreSQL~\cite{postgresql} & 17.11 & Motor del Datawarehouse (esquemas \bd{stg} y \bd{dwh}) \\
                   & pgAdmin 4~\cite{pgadmin} & 9.18.0 & Administración web de la base, Query Tool y ERD Tool \\
                   & Docker y Docker Compose~\cite{docker_compose} & 26.1.5 y 2.26.1 & Contenedores de los servicios \\
    \midrule
    Explotación & Cube (Cube Core)~\cite{cube} & 1.8.19 & Cubo semántico: dimensiones y métricas definidas en YAML \\
                & Streamlit~\cite{streamlit} y Plotly & 1.66.1 y 7.1.0 & Dashboard de visualización web \\
    \midrule
    Calidad y documentación & pytest y ruff & 9.2.0 y 0.17.3 & Pruebas automáticas y revisión de estilo del código \\
                & Playwright & 1.64.1 & Capturas automatizadas de pgAdmin para el informe \\
                & Git y GitHub & --- & Control de versiones del proyecto \\
                & LaTeX (LuaLaTeX) & TeX Live 2026 & Redacción de este informe \\
    \bottomrule
  \end{tabularx}
  \nota{Versiones instaladas y verificadas en el equipo de desarrollo, salvo
  Apache Hop y Cube, que se instalan en el Avance N°2 (ETL) y el Avance N°3
  (cubo). Para ellas se fijaron las últimas versiones estables publicadas al
  5 de noviembre de 2026: Apache Hop 2.20.0 (\bd{downloads.apache.org}) y la
  imagen \bd{cubejs/cube:v1.8.19} (Docker Hub).}
\end{table}

\subsection{Capas de Almacenamiento}

Los datos pasan por cuatro lugares de almacenamiento antes de llegar al
usuario:

\begin{enumerate}[leftmargin=*]
  \item \textbf{\bd{data/raw/}:} archivos originales descargados de cada
        fuente, sin modificar. Permiten repetir el ETL desde cero.
  \item \textbf{\bd{data/processed/}:} archivos intermedios, como la muestra
        de viajes de taxi o los clientes sintéticos generados.
  \item \textbf{Esquema \bd{stg} (staging):} tablas con los datos de cada
        fuente tal como llegan, donde el ETL los limpia y normaliza. Se agrega
        en el Avance N°2.
  \item \textbf{Esquema \bd{dwh}:} el modelo estrella definitivo
        (Sección~\ref{sec:modelo}), que solo recibe datos validados. Está
        implementado en este avance.
\end{enumerate}

Separar \bd{stg} de \bd{dwh} evita que un error de transformación deje datos
inconsistentes en el modelo: si una carga falla, el esquema \bd{dwh} no
cambia.

\subsection{Flujo de Integración (ETL)}

La Tabla~\ref{tab:flujo_etl} resume las etapas del proceso de integración que
se construirá en el Avance N°2. El número de registros leídos, cargados y
rechazados en cada etapa formará la sección Estadísticas de Procesamiento.

\begin{table}[H]
  \centering
  \caption{Etapas del flujo de integración}
  \label{tab:flujo_etl}
  \small
  \begin{tabularx}{\linewidth}{@{}c>{\raggedright\arraybackslash}p{3.2cm}
                                  >{\raggedright\arraybackslash}p{2.4cm}
                                  >{\raggedright\arraybackslash}X@{}}
    \toprule
    \textbf{N°} & \textbf{Etapa} & \textbf{Herramienta} & \textbf{Descripción} \\
    \midrule
    1 & Extracción & Python & Descarga de las fuentes a \bd{data/raw/} y
      muestreo de las fuentes de gran volumen \\
    2 & Clientes sintéticos & Python & Generación de los clientes residentes
      en Chile con las distribuciones del Censo 2024 \\
    3 & Carga a staging & Apache Hop & Lectura de los archivos y carga en el
      esquema \bd{stg} \\
    4 & Transformación & Apache Hop & Limpieza, conversión de unidades y
      moneda, traslado de fechas a 2023--2025, normalización de zonas,
      canales y medios de pago, y asignación de cada transacción a un
      cliente \\
    5 & Carga de dimensiones & Apache Hop & Inserción de los registros nuevos
      en las seis dimensiones, buscando por su llave natural
      (\bd{codigo_<entidad>}) \\
    6 & Carga de hechos & Apache Hop & Búsqueda de las llaves sustitutas de
      cada dimensión e inserción en \bd{fact_venta} \\
    7 & Validación & Apache Hop y SQL & Conteo de registros por etapa,
      registros rechazados y comparación de totales con las fuentes \\
    \bottomrule
  \end{tabularx}
\end{table}

\subsection{Entorno de Ejecución}

El archivo \bd{docker-compose.yml} define los servicios de la solución con
imágenes de versión fija: PostgreSQL (\bd{postgres:17.11}) y pgAdmin
(\bd{dpage/pgadmin4:9.18.0}); Cube (\bd{cubejs/cube:v1.8.19}) se agregará en
el Avance N°3. Los puertos se
publican solo en la interfaz local (\bd{127.0.0.1}), los datos de cada servicio
se guardan en volúmenes de Docker, y las credenciales se leen desde un archivo
\bd{.env} que no se versiona. Con estas decisiones, cualquier integrante
levanta el mismo entorno con un solo comando, en Linux o en Windows (Anexo~B,
Sección~\ref{sec:anexo_entorno}).

El código del proyecto se versiona con Git en un repositorio público de
GitHub (Anexo~C, Sección~\ref{sec:anexo_repositorio}). Los scripts de Python
se revisan con ruff y el Datawarehouse se valida con pruebas automáticas en
pytest (Sección~\ref{sec:verificacion}).
